% 素の `\begin{verbatim}`（LaTeX 標準）。**本文と同じ大きさ**の `\ttfamily` で組まれる。
%
% 逐語ブロックの入口を「本文の文字高と 5% 以上違う」で決めていた頃、`listings` は
% `\ttfamily\small` で組むのがふつうなので通ったが、`verbatim` は字高の差がゼロなので
% 塊として見えず、コードが `body` と `formula` に一行おきに割れていた。
%
% 直す前の実測（本文 4 文が 4 文のままだが、うち 1 文が原文に無い）:
%
%   `def sparse_attention(q, k, v, pattern): scores = scores / math.sqrt(q.shape[-1])
%    … return out We measured the runtime of this implementation on a single
%    accelerator and report the numbers below.`
%
% **直後の本文段落までコードと 1 文になる**ので、訳させると原文に無い文が返る。
% さらにコード 2 行だけが別に `formula` へ落ちて、コードとしても歯抜けになっていた。
% 申告は 3 つとも low。
%
% 正解の取り出し方は他の real-papers/*.tex と同じ（本文 1 文 = \s{...} を 1 行）。
\documentclass[10pt]{article}

\usepackage{mathptmx}
\usepackage[margin=2cm]{geometry}

\newcommand{\s}[1]{#1}

\title{Notes on the Sparse Attention Kernel}
\author{Anonymous}
\date{}

\begin{document}
\maketitle

\section{Implementation}

\s{The kernel is implemented in about thirty lines of Python, reproduced below.}
\s{Line seven is the only place where the sparsity pattern enters the computation.}

\begin{verbatim}
def sparse_attention(q, k, v, pattern):
    scores = q @ k.transpose(-1, -2)
    scores = scores / math.sqrt(q.shape[-1])
    mask = build_mask(pattern, scores.shape)
    scores = scores.masked_fill(mask, float("-inf"))
    weights = torch.softmax(scores, dim=-1)
    weights = apply_dropout(weights, p=0.1)
    out = weights @ v
    out = out.reshape(q.shape[0], -1)
    out = out_proj(out)
    out = layer_norm(out)
    out = residual(out, q)
    return out
\end{verbatim}

\s{We measured the runtime of this implementation on a single accelerator and report the numbers below.}

\section{Results}

\s{The sparse kernel reaches the same accuracy as the dense baseline on every task that we tried.}
\s{Throughput improves by a factor of three at the longest sequence length in the suite.}

\end{document}
